\section{问题三的模型建立与求解}
\label{sec:question-three}

\subsection{问题分析与模型输入}
问题二确定了每张图在各核数下的分核与核内顺序；问题三增加共享只读Cache，关键变化是同一份张量在不同核心上的读取时刻和读取来源。为分开观察硬件变化与调度调整，本章以同图同核数下问题二的固定核数方案$X_B^*$为起点（即该核数下选出的计划，未经继承），包括操作到子图的映射、子图到核心的映射和各核内顺序。先在C配置下评价这份计划，再围绕共享输入产生新候选，由此得到$F_B(X_B^*)$、$F_C(X_B^*)$和最终的$F_C(X_C^*)$。

共享只读Cache容量为1048576 B，读取带宽为250 B/cycle；DDR带宽仍为60 B/cycle，两个读取池独立。Cache初始为空，按先进先出淘汰。逻辑张量标识用于命中查询，实际COPY事件来自问题二的完整核内展开，Spill产生的物理版本也由评估程序处理。每组配对都核对三次评价使用的计划内容是否一致。
\begin{table}[H]\centering\caption{问题三的固定输入和配对量}\label{tab:c-inputs}
\begin{tabularx}{\linewidth}{lY}\toprule
对象 & 值或判定条件\\\midrule
同图同核起点 & 问题二固定核数方案$X_B^*$；C的初始计划与之相同\\
共享Cache & 容量1048576 B；读带宽250 B/cycle；初始队列为空\\
外部存储 & DDR带宽60 B/cycle；写出与未命中读仍走DDR\\
正式响应 & 总工期、结构与恢复搬运、命中和未命中字节、完整计划\\
选择预算 & 常规额度为$\max\{\min(q,2),m+1\}$；$m$为去重后必评锚点数\\\bottomrule
\end{tabularx}\end{table}

\subsection{基于查询与完成事件的FIFO模型}
设$\mathcal C(\tau)$为时刻$\tau$按插入次序排列的逻辑张量队列。合格读入事件$e$在发射时刻$q_e$查询$t_e$，其读取来源由
\begin{equation}
H_e=\mathbf 1\{t_e\in\mathcal C(q_e)\},\qquad
B_e=\begin{cases}250,&H_e=1,\\60,&H_e=0\end{cases}\ \mathrm{B/cycle}
\label{eq:c-query}
\end{equation}
确定。命中读进入Cache池，未命中读进入DDR池；命中只改变本次读取所用的服务池，不把条目移到队尾，查询也不改变FIFO次序。在该COPY\_IN完成时，评估器再尝试把张量插入队尾：若当前仍在队列中则不改变次序；若已经被其他完成事件淘汰，则可能重新插入；若大小超过容量则不插入；空间不足时先逐项淘汰队首。同一时刻的完成、插入和新发射按评估程序的固定事件次序处理，尚未完成的在途读不会被后续查询命中。式\eqref{eq:c-query}中的$B_e$是该次读取所用的服务带宽。

考虑容量为$2b$、三个张量$x,y,z$大小均为$b$的符号例。已有队列$[x,y]$，对$x$的命中不会将其移到队尾。随后$z$完成并插入，队列变为$[y,z]$；如果先前对$x$的在途读取此时才完成，它又能在队尾重新插入，队列变为$[z,x]$。同一组张量的访问次数并不足以确定这段队列演化，完成事件的先后与查询先后都必须考虑。

更关键的是并发查询。两个核心在首个读入完成前向空Cache发射相同张量的请求，两次都未命中；只有后发请求在首次完成插入之后查询且条目未被淘汰，才可能命中。记合格读入总命中字节为$Q_L$，未命中字节为$Q_D$，则报告的字节命中率为
\begin{equation}
h_{\rm byte}=\frac{Q_L}{Q_L+Q_D}.
\label{eq:hit-rate}
\end{equation}
分母为零时取零。该指标只描述读取来源，额外搬运仍按实际生成的COPY事件统计。

DDR和Cache两个共享服务池各自按活动请求分配参考工作量。若池$p$中同时有$m_p(\tau)$个未完成请求，请求$e$的剩余参考工作量$R_e$满足
\begin{equation}
\frac{\mathrm dR_e}{\mathrm d\tau}=-\frac{1}{m_p(\tau)},\qquad
R_e(q_e)=\max\left\{1,\left\lceil\frac{b_e}{B_e}\right\rceil\right\}.
\label{eq:c-service}
\end{equation}
一次命中既可能减轻DDR池竞争，也可能改变后继操作的发射时刻，从而改变关键路径、另一池负载及后续FIFO状态。因此即使$h_{\rm byte}>0$，总工期仍可能不变。

\begin{algorithm}[H]\caption{评估程序中只读FIFO Cache的事件规则}\label{alg:c-fifo}
\begin{algorithmic}[1]
\REQUIRE 已展开的核级执行图、COPY事件、Cache容量及两类带宽
\ENSURE 完成时间线、Cache事件、命中与未命中字节
\STATE 初始化空FIFO队列、DDR池与Cache读池
\WHILE{仍有操作未完成}
\STATE 按评估程序的事件次序处理当前完成事件，推进相应池的剩余工作量
\FOR{当前完成的每个合格COPY\_IN}
\IF{逻辑张量不在队列且大小不超过容量}
\STATE 从队首淘汰至空间充足，并插入该张量
\ENDIF
\ENDFOR
\STATE 更新依赖、流水线与跨核释放就绪状态
\FOR{当前可发射的合格COPY\_IN}
\STATE 在发射时查询队列，选择Cache读池或DDR池并累计字节
\ENDFOR
\STATE 重排两池的预计完成事件并推进事件时间
\ENDWHILE
\end{algorithmic}\end{algorithm}
\subsection{Cache查询时序的解析算例}
为把式\eqref{eq:c-query}与事件先后联系起来，构造一个不参与100图统计的解析小例：三个核心各需读同一份60000 B张量$x$，Cache容量足以容纳$x$。独占DDR读取为$60000/60=1000$ cycle，独占Cache读取为$60000/250=240$ cycle。若三核在时刻0同时发射，三次查询都看到空队列，三个DDR请求平分服务，读取均在时刻3000完成，命中字节为0。若第一核在时刻0发射并于1000完成插入，其余两核在时刻1000才发射且$x$未被淘汰，则两次都命中；两个Cache请求平分读池，各需480 cycle，读取于1480完成。表\ref{tab:c-query-example}把发射、插入和完成三个状态明确分开。
\begin{table}[H]\centering\caption{三核重复读取60000 B张量的解析时序}\label{tab:c-query-example}
\begin{tabularx}{\linewidth}{lYrrr}\toprule
情形 & 发射与插入顺序 & DDR读取 & Cache读取 & 最后一次读完成/cycle\\\midrule
同时查询 & 三核均在0发射，首次插入发生在完成时 & 3次 & 0次 & 3000\\
错峰查询 & 第一核0发射、1000插入，其余两核1000发射 & 1次 & 2次 & 1480\\\bottomrule
\end{tabularx}\end{table}
\begin{figure}[H]\centering
\begin{tikzpicture}[x=1.65cm,y=.9cm,>=Stealth,every node/.style={font=\small}]
\draw[->] (0,0)--(3.5,0) node[right] {$\tau/1000$ cycle};
\foreach \x/\lab in {0/0,1/1,1.48/1.48,3/3}{\draw (\x,.06)--(\x,-.06) node[below] {\lab};}
\node[left] at (0,1.0) {错峰查询};
\node[left] at (0,2.4) {同时查询};
\draw[blue!65!black,very thick] (0,2.4)--(3,2.4);
\node[above] at (1.5,2.4) {3次DDR未命中并发};
\draw[blue!65!black,very thick] (0,1.0)--(1,1.0);
\node[above] at (.5,1.0) {1次DDR};
\draw[green!55!black,very thick] (1,1.0)--(1.48,1.0);
\node[anchor=west] at (1.58,1.0) {2次Cache命中};
\draw[dashed] (1,.12)--(1,2.9);
\node[anchor=west] at (1.65,3.25) {首读完成并插入：$\tau=1000$};
\end{tikzpicture}
\caption{相同访问次数下的两种查询顺序；仅为解析示例}\label{fig:c-query-timeline}
\end{figure}
错峰发射在真实计划中须由前置计算、依赖或队列中的其他操作形成，这些操作本身也占用资源。

\subsection{共享输入方案与求解}
\textbf{Step 1：生成候选。}以$X_B^*$为种子，尝试迁移、同核合并、核内调序以及跨核分散，并查找跨核重复读取的逻辑张量。对按张量大小和潜在重复读量排序的有限张量集合，枚举相关消费者组或节点在现有核心序列中的前移、后移和跨核插入位置；对满足容量条件的共享输入，再指定一个领读核，将其消费者移至该核序列前部，并把其他消费者后移到各自序列尾部，形成错峰查询候选。候选只保留合法映射与顺序；共享张量集合随当前Cache容量过滤，因而容量扰动会同步改变候选组合。不合法的构造不进入完整评价。

\textbf{Step 2：轻量排序。}按路径递推、最大核心计算量、结构COPY的DDR参考工作量和容量压力估计候选代价；C另外在收缩后的商图上回放逻辑COPY查询的相对时序，记录查询时的FIFO命中代理、未命中量和潜在Cache服务工作量。该回放保留每核顺序、跨核释放和完成后插入规则，但不模拟服务池的并发；正式工期、命中和淘汰由问题三的事件模拟给出。

\textbf{Step 3：完整评价与选优。}把B最终计划$X_B^*$、整图、分量装桶和低核数继承计划中去重后的方案置于必评锚点，其余C候选按轻量分数排列。设去重后必评数为$m$、请求额度为$q$，常规完整调用额度为$\max\{\min(q,2),m+1\}$。展开失败只记录错误，不占完整调用额度；完整评价失败占用一次额度。首个成功结果记为初始方案，最终按工期、额外搬运的字典序在全部成功结果中选定$X_C^*$。得到首个C结果后，从其Cache事件中提取重复未命中张量，生成事件锚定的合法调序候选；若仍有额度，再按候选代表补评。下文500组B/C配对均以同一$X_B^*$为C起点，用于同计划效应分解，与继承方案的总体曲线分开报告。

\begin{algorithm}[H]\caption{问题三候选选择与同计划对照}\label{alg:c-select}
\begin{algorithmic}[1]
\REQUIRE 问题二固定核数方案$X_B^*$、图、核数、固定硬件参数及请求完整调用额度$q$
\ENSURE 问题三选中计划$X_C^*$及三段工期比
\STATE 保留$X_B^*$，生成迁移、同核合并、调序、分散及领读核错峰代表
\STATE 逐个检查合法性、按完整计划去重，计算轻量排序分数
\STATE 将$X_B^*$、整图、分量装桶和低核数继承计划置于去重后的必评队列
\STATE 令$m$为必评数，完整调用额度为$\max\{\min(q,2),m+1\}$
\WHILE{未达到常规完整调用额度且还有未尝试候选}
\STATE 取必评队列或轻量排序最靠前的未尝试候选，先展开检查再做C场景事件模拟
\STATE 展开失败不占额度；完整评价保存成功结果或完整错误记录
\ENDWHILE
\STATE 从首个成功结果提取重复未命中事件，生成事件锚定调序候选并在剩余额度内补评
\STATE 在成功结果中按工期、额外搬运的字典序选$X_C^*$
\STATE 核对C初始计划与$X_B^*$一致，计算式\eqref{eq:cache-factors}
\end{algorithmic}\end{algorithm}

\subsection{共享输入聚合与错峰的解析边界}
仍用60000 B共享输入，令五个V流水线操作各计算1000 cycle，暂不计输出。下面比较两种给定的局部请求时序下服务池的边界工期。全部聚到一核，只读一次DDR却顺序完成五个计算，理论完成时间为$1000+5\times1000=6000$ cycle。五核分散且同时查询空Cache时，五次未命中读平分DDR，读阶段持续$5\times60000/60=5000$ cycle，再各算1000 cycle，总时间同为6000。若只有第一核在0查询，其余四核由外部前置关系约束到1000才查询，四次命中合计240000 B，Cache读阶段需$240000/250=960$ cycle，后四核在2960完成。读取时序与Cache服务能力共同决定局部读计算完成时间。
\begin{table}[H]\centering\caption{五个V操作共享60000 B输入的解析边界（非正式实验）}\label{tab:c-five-example}
\begin{tabularx}{\linewidth}{lYr}\toprule
结构 & 读取与执行条件 & 解析工期/cycle\\\midrule
单核聚合 & 一次DDR读取，五次V操作串行 & 6000\\
五核分散、无Cache & 五次DDR读同时开始、共享60 B/cycle & 6000\\
五核分散、Cache同时查询 & 空队列下五次均未命中 & 6000\\
五核分散、其余四核错峰 & 后四次命中，Cache带宽250 B/cycle & 2960\\
同上、带宽减半 & 后四次命中，Cache带宽125 B/cycle & 3920\\
同上、带宽加倍 & 后四次命中，Cache带宽500 B/cycle & 2480\\
同上、容量不足60000 B & 张量不能插入，后四次仍未命中 & 6000\\\bottomrule
\end{tabularx}\end{table}
五核分散在250 B/cycle条件下，第一核于2000完成计算；其余四核从1000起共读960 cycle，再计算1000 cycle，于2960完成。带宽减半和加倍时，将960分别改为1920、480，即得3920、2480。

\subsection{同计划对照与调度适配}
同一图同一核数取得三个工期：$F_B(X_B^*)$、$F_C(X_B^*)$、$F_C(X_C^*)$。定义
\begin{equation}
\begin{split}
S_{\rm hw}&=\frac{F_B(X_B^*)}{F_C(X_B^*)},\\
S_{\rm adapt}&=\frac{F_C(X_B^*)}{F_C(X_C^*)},\\
S_{\rm opt}&=\frac{F_B(X_B^*)}{F_C(X_C^*)}=S_{\rm hw}S_{\rm adapt}.
\end{split}
\label{eq:cache-factors}
\end{equation}
第一个比值只改硬件场景，第二个只在C中换计划，最后一个是B最终到C最终的综合工期比。500组配对的乘积恒等式最大数值残差为$2.22\times10^{-16}$。表\ref{tab:cache-factors}每列分别对100张图的逐例比值取算术平均，因此列均值不满足乘积恒等式。
\begin{table}[H]\centering\caption{每核100图的Cache收益三段分解}\label{tab:cache-factors}
\begin{tabular}{crrrr}\toprule
核数 & 硬件效应均值 & 适配效应均值 & 综合效应均值 & 适配改善图数\\\midrule
1 & 1.0044 & 1.0249 & 1.0293 & 31\\
2 & 1.0068 & 1.0000 & 1.0068 & 1\\
3 & 1.0088 & 1.0011 & 1.0099 & 1\\
4 & 1.0138 & 1.0011 & 1.0150 & 7\\
5 & 1.0125 & 1.0011 & 1.0136 & 7\\\bottomrule
\end{tabular}\end{table}
图\ref{fig:cache-direct-comparison}给出$F_B(X_B^*)$、$F_C(X_B^*)$和$F_C(X_C^*)$的绝对平均工期及后二者相对$F_B(X_B^*)$的逐图比值均值；绝对工期受图规模影响，比值均值才反映硬件效应与调度适配的相对幅度。
\samplefigure{cache_direct_comparison}{问题三无Cache、固定计划Cache与重调度Cache的同核直接比较}{fig:cache-direct-comparison}

图\ref{fig:cache-coverage-radar}给出每核100图中工期严格缩短的占比：五核硬件效应、C内适配与综合改善分别覆盖54、7、58图；单核有31图受益于C内适配，二、三核各只有1图。硬件与适配的改善图有3图重合。
\samplefigure{cache_effect_coverage_radar}{五个核数下硬件、适配与综合工期改善的逐图覆盖率分组条形图}{fig:cache-coverage-radar}

五核100图中，$S_{\rm hw}>1$有54图，恰为1有46图，没有退化图；$C(X_B^*)$的字节加权命中率为16.79\%。54张改善图的工期降幅中位数为0.4428\%，均值为2.0074\%（折合速度增幅均值2.31\%），其中36图降幅不足1\%。
\samplefigure{cache_hw_effect_pie}{五核硬件效应的图数构成及改善幅度分类条形图}{fig:cache-hardware-pie}

对\Case{080}，固定B计划的工期由113455降至90408 cycle，$S_{\rm hw}=1.2549$，工期降幅20.3138\%；对\Case{055}，$C(X_B^*)$存在189952字节命中，但工期两侧均为25635 cycle——命中未落在决定尾部完成时刻的关键路径上。图\ref{fig:cache-hit-gain}给出五核100图的$C(X_B^*)$命中率与同计划工期改善，零收益点也保留在图中。
\samplefigure{cache_hit_gain_scatter}{五核字节命中率与同计划工期改善的逐图关系}{fig:cache-hit-gain}

\begin{table}[H]\centering\caption{四张代表图的五核三段工期与Cache访问}\label{tab:c-cases}
\begin{tabular}{lrrrrr}\toprule
图 & $F_B(X_B^*)$ & $F_C(X_B^*)$ & $F_C(X_C^*)$ & $C(X_B^*)$命中/B & $S_{\rm hw}$\\\midrule
\Case{026} & 45246 & 44374 & 41655 & 21936 & 1.0197\\
\Case{053} & 1472686 & 1428109 & 1428109 & 5155968 & 1.0312\\
\Case{055} & 25635 & 25635 & 25635 & 189952 & 1.0000\\
\Case{080} & 113455 & 90408 & 90408 & 2228224 & 1.2549\\\bottomrule
\end{tabular}\end{table}
表\ref{tab:c-cases}把影响分成三种可见情形：\Case{055}读来源有变化但完成时间不变；\Case{080}固定计划的硬件效应较大；\Case{026}则在硬件效应之外还有C内调度适配。\Case{053}在四图中命中字节最多，硬件效应却仅1.0312，命中量与时间收益并不成比例。

把同计划比较扩展到全部500个图-核组合，图\ref{fig:cache-hit-gain-3d}按核数分面，在二维坐标中表示$C(X_B^*)$字节命中率与$100[1-F_C(X_B^*)/F_B(X_B^*)]$工期降幅。205组缩短了工期，另有141组虽发生命中但工期不变；154组既无命中也无工期变化，在各分面原点叠印。这500组中，所有同计划工期缩短的结果都伴随非零命中，但有命中不等于工期下降。
\samplefigure{cache_hit_gain_3d}{500组同计划配对的分核命中率与工期降幅分面散点图}{fig:cache-hit-gain-3d}

\Case{026}则说明第二段效应：同一B计划在C下的起点为44374 cycle，选中重复读取相关候选后为41655 cycle，$S_{\rm adapt}=1.0653$。五核仅7图出现适配改善，这类收益集中在少数计划。500组总计有205组硬件效应改善、295组持平；47组适配改善、453组持平；综合改善232组、持平268组，三类比值的中位数均为1。

进一步检查两段收益是否发生在同一批图。以$S_{\rm hw}>1$和$S_{\rm adapt}>1$对五核100图交叉分类，等于1记为未改善。表\ref{tab:c-effect-cross}的综合工期降幅先按每图的$100[1-F_C(X_C^*)/F_B(X_B^*)]$计算，再在组内等权平均。
\begin{table}[H]\centering\small
\caption{五核100图的硬件效应与调度适配交叉分布}
\label{tab:c-effect-cross}
\begin{tabular}{lrr}\toprule
效应分类 & 图数 & 综合工期降幅均值/\%\\\midrule
仅同计划硬件效应 & 51 & 2.0778\\
仅C内调度适配 & 4 & 1.1170\\
两种效应兼有 & 3 & 2.9412\\
两种效应均无 & 42 & 0.0000\\\bottomrule
\end{tabular}
\end{table}
硬件效应改善的54图与适配改善的7图有3图重合，因此综合改善为58图。\Case{040}的B最终与C同计划工期均为110337 cycle，更换C计划后降至110105 cycle，最终计划的命中字节为0。这232 cycle来自计划之间的差异而非命中；另有51张图在保留B计划时就已缩短工期。交叉分类表明，是否更换计划与是否受益于Cache硬件是两个独立的判断步骤。

\subsection{问题三结果}
问题三采用与前两问相同的继承方案：对目标核数$K$，在已评价的$1,\ldots,K$核方案及整图单核基准中取Makespan最短者，低核方案以空核列表补齐。相对统一整图单核基准的五核逐例加速比均值为3.4272、中位数4.0719，所有图均不慢于基准。A、B单核曲线按定义显示为1，C单核保留真实工期比。
\begin{table}[H]\centering\caption{问题三：100图逐核加速比（继承方案；慢例按未舍入工期严格比较）}\label{tab:c-results}
\begin{tabular}{crrrr}\toprule
核数 & 平均加速比 & 中位数 & 慢于整图单核图数 & 首个成功至最终平均降幅\\\midrule
1 & 1.0468 & 1.0000 & 0 & 2.1873\%\\
2 & 1.7232 & 1.9703 & 0 & 0.0215\%\\
3 & 2.3335 & 2.8340 & 0 & 0.0215\%\\
4 & 2.9496 & 3.6732 & 0 & 0.0215\%\\
5 & 3.4272 & 4.0719 & 0 & 0.0828\%\\\bottomrule
\end{tabular}\end{table}
\samplefigure{speedup_C}{问题三加速比均值、中位数及均值95\%重采样区间}{fig:speedup-c}

图\ref{fig:speedup-c}的浅色带为100张图重采样1500次得到的均值95\%区间。均值由4核的2.9496升至5核的3.4272，增加0.4776；由4核到5核有69张图工期缩短、31张图保持不变，没有图因增加一个核心而变慢。

\subsection{问题三灵敏度分析：L2容量与带宽}
固定六张代表图的五核B最终计划，分别把L2 Cache容量与读带宽减半、加倍，考察硬件参数对命中、工期和切图调度的影响。六图按五核B计划的基线工期、初始命中率和规模分层，覆盖小工期低命中图（012、081）、中等工期高命中图（051）、容量边界图（062）以及大工期图（016、067）；样本在实验前固定，未按扰动后的结果挑选。“固定计划”只改变硬件参数，“重优化”则重新生成C候选。24个固定计划结果全部成功；重优化候选中36条成功、12条因依赖环检查失败，失败候选不计入均值。24组重优化最终全部选回固定B计划。
\paragraph{参数作用的直接判断。}容量减半使六图平均工期增加$0.2701\%$，容量加倍使其减少$0.0517\%$；变化集中在容量边界附近的图\Case{062}，其工期分别增加$1.6206\%$和减少$0.3103\%$。带宽减半和加倍的平均工期变化为$-0.0012\%$和$+0.0089\%$，方向未呈单调关系。四种扰动下重优化均回到同一份固定计划，说明在这组数据中，L2参数主要改变读请求的服务时序，还没有带来足以改变切图、分核或核内顺序的收益。
\begin{table}[H]\centering\small\caption{六图四种Cache扰动的绝对工期均值与逐图命中率（固定计划与重优化）}\label{tab:cache-sensitivity}
\begin{tabular}{lrrrrr}\toprule
配置 & \shortstack{固定绝对工期均值\\/cycle} & \shortstack{重优化绝对工期均值\\/cycle} & \shortstack{逐图命中率\\均值} & \shortstack{逐图相对默认\\变化均值} & 失败候选\\\midrule
容量减半 & 4343744.7 & 4343744.7 & 10.71\% & +0.2701\% & 3\\
容量加倍 & 4338303.3 & 4338303.3 & 14.73\% & -0.0517\% & 3\\
带宽减半 & 4339193.8 & 4339193.8 & 14.63\% & -0.0012\% & 3\\
带宽加倍 & 4339315.0 & 4339315.0 & 14.63\% & +0.0089\% & 3\\\bottomrule
\end{tabular}\end{table}
表中两列工期均值是六张图绝对周期的算术平均；“逐图相对默认变化均值”先按每张图计算$100(F_{\rm variant}/F_{\rm default}-1)$，再在六图间等权平均。“命中率均值”也先逐图计算再平均；若汇总六图的命中与访问字节，加权命中率依次为0.7876\%、1.6757\%、1.5597\%、1.5597\%。

表\ref{tab:c-sensitivity-cases}逐图给出$100(F_{\rm variant}/F_{\rm default}-1)$，正值表示变慢。容量减半与加倍时只有\Case{062}的工期变化；带宽减半影响两图，带宽加倍影响三图，幅度都很小且方向不一。容量的影响集中在驻留区间接近容量边界的图。由于只考察了六张图和有限的候选，本文只把结论限定为“在当前候选范围内未观察到计划改变”，不推断L2参数对切图没有作用。
\begin{table}[H]\centering\small\caption{六张代表图固定计划的Cache参数逐例响应（\%）}\label{tab:c-sensitivity-cases}
\begin{tabular}{lrrrrr}\toprule
图 & 默认工期/cycle & 容量减半 & 容量加倍 & 带宽减半 & 带宽加倍\\\midrule
\Case{012} & 13406 & 0.0000 & 0.0000 & 0.0000 & 0.0000\\
\Case{016} & 7740666 & 0.0000 & 0.0000 & 0.0000 & +0.0006\\
\Case{051} & 627502 & 0.0000 & 0.0000 & -0.0209 & +0.0105\\
\Case{062} & 1690862 & +1.6206 & -0.3103 & +0.0135 & +0.0420\\
\Case{067} & 15931687 & 0.0000 & 0.0000 & 0.0000 & 0.0000\\
\Case{081} & 30943 & 0.0000 & 0.0000 & 0.0000 & 0.0000\\\bottomrule
\end{tabular}\end{table}

命中率与工期并不同步。\Case{051}容量减半后字节命中率从42.39\%降至29.35\%，工期仍为627502 cycle；改变Cache读带宽时，命中率均为42.39\%，工期分别比默认短0.0209\%和长0.0105\%。\Case{062}容量减半使工期增加27402 cycle，容量加倍使工期减少5246 cycle，对应命中率为1.43\%和4.97\%。这些结果说明，命中增加能否缩短尾部完成时间，还取决于读事件所处执行路径及共享池的服务顺序。

\paragraph{小结。}
五核固定计划中54图因Cache缩短工期、46图不变，调整C计划只在7图上进一步改善；六图参数扰动的工期影响集中在容量边界附近的少数图。
